\documentclass[public]{nervtech-overview}

\nervsetup{
  company={NervTech},
  project={Node v0},
  document={Initial Project Overview},
  subtitle={A plain-language briefing for stakeholders, collaborators, and portfolio reviewers},
  docid={NT-N-OVR-001},
  revision={A},
  status={Draft Brief},
  author={Anhadh Virk},
  owner={NervTech Systems Engineering},
  date={2026-06-25},
  audience={Non-technical stakeholders; Portfolio reviewers},
  tagline={Node: The Building Block of NervTech Solutions},
  purpose={To pave the way for NervTech robotic systems, a robust, multipurpose actuator must be developed that has high precision, bandwith and torque capabilities. The Node v0 fills this need, while being 
           developed in house without the extra integration time and cost incurred by OEM integrated actuators.},
  summary={This overview is intentionally non-technical. It covers the current stakeholder need for the product, how it fits into the growth of NervTech, the requirements that define success for the product and finally
          a roadmap to achieve this success. It avoids equations and low-level implementation detail so the project can be understood by reviewers, collaborators, mentors, and potential employers.}
}

\begin{document}
\maketitle
\tableofcontents
\newpage

\section{Project Snapshot}

\projectsnapshot
  {A compact robotic actuator module that combines motion generation, sensing, electronics, and control in one reusable unit.}
  {Robots become easier to build when the joint actuator is reliable, measurable, and simple to integrate.}
  {The project is in early architecture and prototype planning.}
  {A bench-tested actuator module with plots, photos, CAD renders, and a short demonstration video.}

\section{What the Project Is}

\begin{overviewbox}[Plain-Language Summary]
The Integrated Robotic Actuator is a self-contained robotic joint module. Instead of treating the motor, gearbox, sensors, electronics, and software as separate parts, the project brings them together into one engineered actuator package.
\end{overviewbox}

The goal is not only to make something move. The goal is to show a complete robotics engineering workflow: research, design, prototyping, testing, documentation, and demonstration.

\section{Why It Matters}

\begin{valuebox}[Project Value]
A well-designed actuator is a building block for many robotic systems, including arms, mobile robots, legged robots, grippers, and educational robotics platforms.
\end{valuebox}

For a portfolio, this project is valuable because it connects mechanical design, electronics, embedded control, safety testing, data collection, and clear communication.

\section{Who This Is For}

\begin{stakeholdertable}
\stakeholderrow{Engineering reviewers}{Whether the project has credible technical depth.}{Shows architecture, testing intent, and visible deliverables.}
\stakeholderrow{Research mentors}{Whether the project can connect to control, learning, or neuromorphic sensing.}{Frames the project as a reusable experimental platform.}
\stakeholderrow{Non-technical stakeholders}{What the project does and why it matters.}{Explains the project without relying on formulas or jargon.}
\end{stakeholdertable}

\section{Current Status}

\begin{statusbox}[Now]
The project is at the point where the major design direction should be selected, documented, and validated through a small bench prototype.
\end{statusbox}

\begin{milestonetable}
\milestonerow{Concept}{Define the actuator purpose and intended demonstration.}{Current}{In progress}
\milestonerow{Architecture}{Choose motor, sensing, gearbox, controller, and package strategy.}{Next}{Planned}
\milestonerow{Prototype}{Build and test a first integrated bench unit.}{Following}{Planned}
\milestonerow{Portfolio}{Publish video, plots, CAD renders, README, and lessons learned.}{Final}{Planned}
\end{milestonetable}

\section{Roadmap}

\begin{roadmaptable}
\roadmaprow{Phase 1}{Research and concept selection.}{Preliminary research note}
\roadmaprow{Phase 2}{Technical design and requirements.}{Engineering design document}
\roadmaprow{Phase 3}{Prototype build and bench testing.}{Photos, plots, test logs}
\roadmaprow{Phase 4}{Portfolio packaging and public explanation.}{GitHub README and demo video}
\end{roadmaptable}

\section{Communication Positioning}

\begin{plainstakeholderbox}[How to Explain It]
NervTech is building a compact actuator module that can become the motion core of future robotic systems. The project demonstrates not just a part, but an engineering process: make a design decision, build it, test it, learn from it, and document the result clearly.
\end{plainstakeholderbox}

\begin{riskbox}[Communication Risk]
Avoid describing the project as a finished product too early. Present it as a serious prototype and learning platform until the actuator has been built and tested.
\end{riskbox}

\end{document}
